\documentclass[10pt, letterpaper]{article}

% Packages:
\usepackage[
    ignoreheadfoot, % set margins without considering header and footer
    top=2 cm, % seperation between body and page edge from the top
    bottom=2 cm, % seperation between body and page edge from the bottom
    left=2 cm, % seperation between body and page edge from the left
    right=2 cm, % seperation between body and page edge from the right
    footskip=1.0 cm, % seperation between body and footer
    % showframe % for debugging 
]{geometry} % for adjusting page geometry
\usepackage{titlesec} % for customizing section titles
\usepackage{tabularx} % for making tables with fixed width columns
\usepackage{array} % tabularx requires this
\usepackage[dvipsnames]{xcolor} % for coloring text
\definecolor{primaryColor}{RGB}{0, 0, 0} % define primary color
\usepackage{enumitem} % for customizing lists
\usepackage{fontawesome5} % for using icons
\usepackage{amsmath} % for math
\usepackage[
    pdftitle={Lucas Falcão's CV},
    pdfauthor={Lucas Falcão},
    pdfcreator={Me},
    colorlinks=true,
    urlcolor=primaryColor
]{hyperref} % for links, metadata and bookmarks
\usepackage[pscoord]{eso-pic} % for floating text on the page
\usepackage{calc} % for calculating lengths
\usepackage{bookmark} % for bookmarks
\usepackage{lastpage} % for getting the total number of pages
\usepackage{changepage} % for one column entries (adjustwidth environment)
\usepackage{paracol} % for two and three column entries
\usepackage{ifthen} % for conditional statements
\usepackage{needspace} % for avoiding page brake right after the section title
\usepackage{iftex} % check if engine is pdflatex, xetex or luatex

%\addbibresource{own-bib.bib}


% Ensure that generate pdf is machine readable/ATS parsable:
\ifPDFTeX
    \input{glyphtounicode}
    \pdfgentounicode=1
    \usepackage[T1]{fontenc}
    \usepackage[utf8]{inputenc}
    \usepackage{lmodern}
\fi

\usepackage{charter}

% Some settings:
\raggedright
\AtBeginEnvironment{adjustwidth}{\partopsep0pt} % remove space before adjustwidth environment
\pagestyle{empty} % no header or footer
\setcounter{secnumdepth}{0} % no section numbering
\setlength{\parindent}{0pt} % no indentation
\setlength{\topskip}{0pt} % no top skip
\setlength{\columnsep}{0.15cm} % set column seperation
\pagenumbering{gobble} % no page numbering

\titleformat{\section}{\needspace{4\baselineskip}\bfseries\large}{}{0pt}{}[\vspace{1pt}\titlerule]

\titlespacing{\section}{
    % left space:
    -1pt
}{
    % top space:
    0.3 cm
}{
    % bottom space:
    0.2 cm
} % section title spacing

\renewcommand\labelitemi{$\vcenter{\hbox{\small$\bullet$}}$} % custom bullet points
\newenvironment{highlights}{
    \begin{itemize}[
        topsep=0.10 cm,
        parsep=0.10 cm,
        partopsep=0pt,
        itemsep=0pt,
        leftmargin=0 cm + 10pt
    ]
}{
    \end{itemize}
} % new environment for highlights


\newenvironment{highlightsforbulletentries}{
    \begin{itemize}[
        topsep=0.10 cm,
        parsep=0.10 cm,
        partopsep=0pt,
        itemsep=0pt,
        leftmargin=10pt
    ]
}{
    \end{itemize}
} % new environment for highlights for bullet entries

\newenvironment{onecolentry}{
    \begin{adjustwidth}{
        0 cm + 0.00001 cm
    }{
        0 cm + 0.00001 cm
    }
}{
    \end{adjustwidth}
} % new environment for one column entries

\newenvironment{twocolentry}[2][]{
    \onecolentry
    \def\secondColumn{#2}
    \setcolumnwidth{\fill, 4.5 cm}
    \begin{paracol}{2}
}{
    \switchcolumn \raggedleft \secondColumn
    \end{paracol}
    \endonecolentry
} % new environment for two column entries

\newenvironment{threecolentry}[3][]{
    \onecolentry
    \def\thirdColumn{#3}
    \setcolumnwidth{, \fill, 4.5 cm}
    \begin{paracol}{3}
    {\raggedright #2} \switchcolumn
}{
    \switchcolumn \raggedleft \thirdColumn
    \end{paracol}
    \endonecolentry
} % new environment for three column entries

\newenvironment{header}{
    \setlength{\topsep}{0pt}\par\kern\topsep\centering\linespread{1.5}
}{
    \par\kern\topsep
} % new environment for the header

\newcommand{\placelastupdatedtext}{% \placetextbox{<horizontal pos>}{<vertical pos>}{<stuff>}
  \AddToShipoutPictureFG*{% Add <stuff> to current page foreground
    \put(
        \LenToUnit{\paperwidth-2 cm-0 cm+0.05cm},
        \LenToUnit{\paperheight-1.0 cm}
    ){\vtop{{\null}\makebox[0pt][c]{
        \small\color{gray}\textit{Last updated in September 2024}\hspace{\widthof{Last updated in September 2024}}
    }}}%
  }%
}%

% save the original href command in a new command:
\let\hrefWithoutArrow\href

% new command for external links:


\begin{document}
    \newcommand{\AND}{\unskip
        \cleaders\copy\ANDbox\hskip\wd\ANDbox
        \ignorespaces
    }
    \newsavebox\ANDbox
    \sbox\ANDbox{$|$}

    \begin{header}
        \fontsize{25 pt}{25 pt}\selectfont Lucas Falcão

        \vspace{5 pt}

        \normalsize
        \mbox{São Paulo, Brazil}%
        \kern 5.0 pt%
        \AND%
        \kern 5.0 pt%
        \mbox{\hrefWithoutArrow{mailto:lucas.falcao@ieps.org.br}{lucas.falcao@ieps.org.br}}%
        \kern 5.0 pt%
        \AND%
        \kern 5.0 pt%
        \mbox{\hrefWithoutArrow{https://www.linkedin.com/in/lucas-falcao9113}{linkedin.com/in/lucas-falcao}}%
        \kern 5.0 pt%
        \AND%
        \kern 5.0 pt%
        \mbox{\hrefWithoutArrow{https://github.com/lfalcao9113}{github.com/lfalcao9113}}%
    \end{header}

    \vspace{5 pt - 0.3 cm}


\section{Educação}

\begin{twocolentry}{
    2021 -- Presente
}
    \textbf{Fundação Getulio Vargas (FGV)}, Doutorando em Administração Pública e Governo
\end{twocolentry}

\vspace{0.10 cm}
\begin{onecolentry}
    \begin{highlights}
        \item Pesquisador visitante na Duke University, Sanford School of Public Policy (2023 -- 2024)
        \item Foco da pesquisa: economia da saúde; meio ambiente; e Amazônia.
    \end{highlights}
\end{onecolentry}

\vspace{0.10 cm}
\begin{twocolentry}{
    2017 -- 2019
}
    \textbf{Universidade Federal do ABC (UFABC)}, Mestrado em Economia
\end{twocolentry}

\vspace{0.10 cm}
\begin{onecolentry}
    \begin{highlights}
        \item Dissertação: “Existe um trade-off entre eficiência e cooperativismo? Evidências de cooperativas de trabalho brasileiras”
    \end{highlights}
\end{onecolentry}

\vspace{0.10 cm}
\begin{twocolentry}{
    2013 -- 2017
}
    \textbf{Universidade Federal do ABC (UFABC)}, Bacharelado em Economia
\end{twocolentry}

\vspace{0.10 cm}
\begin{twocolentry}{
    2013 -- 2016
}
    \textbf{Universidade Federal do ABC (UFABC)}, Bacharelado em Ciências e Humanidades
\end{twocolentry}


\section{Experiência}

\begin{twocolentry}{
    2025 -- Presente
}
    \textbf{Pesquisador}, Instituto de Estudos para Políticas de Saúde (IEPS)
\end{twocolentry}

\vspace{0.10 cm}
\begin{onecolentry}
    \begin{highlights}
        \item Foco: desigualdades raciais em saúde no Brasil.
    \end{highlights}
\end{onecolentry}

\vspace{0.2 cm}

\begin{twocolentry}{
    2021 -- 2024
}
    \textbf{Assistente de Pesquisa}, Instituto de Estudos para Políticas de Saúde (IEPS) e Amazônia 2030
\end{twocolentry}

\vspace{0.10 cm}
\begin{onecolentry}
    \begin{highlights}
        \item Projeto: “Saúde na Amazônia Legal: Evolução Recente e Desafios em Perspectiva Comparada.”
    \end{highlights}
\end{onecolentry}

\vspace{0.2 cm}

\begin{twocolentry}{
    2019 -- 2020
}
    \textbf{Assistente de Pesquisa}, Centro de Política e Economia do Setor Público (CEPESP)
\end{twocolentry}

\vspace{0.10 cm}
\begin{onecolentry}
    \begin{highlights}
        \item Projeto: “Investigação sobre a atuação política de empresas do setor de saúde no Brasil.”
    \end{highlights}
\end{onecolentry}

\vspace{0.2 cm}

\begin{twocolentry}{
    2019 -- 2020
}
    \textbf{Pesquisador}, Centro Brasileiro de Análise e Planejamento (Cebrap)
\end{twocolentry}

\vspace{0.10 cm}
\begin{onecolentry}
    \begin{highlights}
        \item Projeto: “Efeitos socioeconômicos da expansão da soja entre 1991 e 2010.”
    \end{highlights}
\end{onecolentry}

\vspace{0.2 cm}

\begin{twocolentry}{
    2019
}
    \textbf{Assistente de Pesquisa}, DevLab, Duke University
\end{twocolentry}

\vspace{0.10 cm}
\begin{onecolentry}
    \begin{highlights}
        \item Projeto: “Clientelismo e Partidos Políticos no Brasil.”
    \end{highlights}
\end{onecolentry}


\section{Publicações}

\input{publications}


\section{Habilidades}

\begin{twocolentry}{
    R, STATA, Python
}
    \textbf{Análise de Dados \& Econometria}
\end{twocolentry}

\vspace{0.10 cm}
\begin{onecolentry}
    \begin{highlights}
        \item Manipulação de microdados, modelagem econométrica e aprendizado de máquina.
    \end{highlights}
\end{onecolentry}

\vspace{0.2 cm}

\begin{twocolentry}{
    QGIS, R (sf, ggmap, Shiny)
}
    \textbf{Análise Geoespacial \& Visualização}
\end{twocolentry}

\vspace{0.10 cm}
\begin{onecolentry}
    \begin{highlights}
        \item Cartografia, análise espacial e desenvolvimento de aplicativos web interativos.
    \end{highlights}
\end{onecolentry}


\section{Premiações}

\begin{twocolentry}{
    2018
}
    \textbf{Prêmio de Excelência Acadêmica}, CECS-UFABC
\end{twocolentry}

\vspace{0.2 cm}

\begin{twocolentry}{
    2017
}
    \textbf{Melhor Monografia}, Bacharelado em Economia (UFABC)
\end{twocolentry}


\end{document}